Hakates looma uut süsteemi, mis kasutab andmete vahetamiseks eDelivery raamistikku,
peab arendaja leidma enda projekti jaoks sobilikud valmiskomponendid,
mis võimaldavad eDelivery sõnumeid osapoolte vahel saata.

eDelivery standard on tarnija- ja platvormineutraalne, mis tähendab et üks kindel pakkuja
ei kontrolli süsteemi arhitektuuri, ning kõigil on võimalik luua oma kohandatud lahendus.
See tähendab seda, et täna on turul väga palju erinevaid eDelivery tooteid ja teenuseid, 
mille vahel saab valikuid teha \cite{ec_edelivery_building_block}. 

Selle lõputöö raames on vaja leida ja kasutusele võtta kolme tüüpi eDelivery lahendusi. Nendeks on:
\begin{enumerate}
  \item eDelivery juurdepääsupunkt (AP)
  \item eDelivery teenuse metaandmete avaldaja (SMP)
  \item eDelivery teenuse metaandmete leidja (SML)
\end{enumerate}

\section{eDelivery juurdepääsupunkt (AP)}
Euroopa komisjoni eDelivery tutvustaval veebilehel on välja toodud abistavaid linke ja juhendeid nendele,
kes soovivad seda kasutada \cite{ec_edelivery_building_block}. Nende linkide seas on viide veebilehele,
mis aitab arendajatel leida ja võrrelda AS4 tarkvara või teenuseid, 
mis on läbinud eDelivery vastavustestimise \cite{ec_edelivery_as4_conformant_solutions}.

Autori eesmärk selle lõputöö raames on luua lahendus, mis töötaks lokaalselt,
ning millel ei ole liigseid väliseid sõltuvusi. 
Selle tõttu on autor otsustanud kasutada ainult eDelivery tooteid, mis on kättesaadavad tasuta, 
ning mille lähtekood on avalik.
See tähendab, et eDelivery teenused ning teenusepakkujaid see uurimustöö ei analüüsi.






